% Section 5: Results
\section{Results}
\label{sec:results}

All five sub-hypotheses passed their gates. We report in RQ order.

\subsection{RQ1: Capability Independently Predicts LC Preference (H-E1)}

After controlling for \texttt{avg\_length}:

\begin{quote}
$r_{\mathrm{partial}} = 0.9851$, $p = 1.69 \times 10^{-170}$ (one-tailed),
Bootstrap 95\% CI $[0.976, 0.988]$, $N=223$
\end{quote}

VIF = 1.764 confirms no multicollinearity (Figure~\ref{fig:vif}). Figure~\ref{fig:scatter}
shows the bivariate scatter ($\rho = 0.966$); Figure~\ref{fig:partial_reg} shows the partial
regression plot ($r_{\mathrm{partial}} = 0.985$).

\textbf{The counterintuitive finding}: partial correlation (0.985) exceeds bivariate (0.94).
Verbosity acts as a mild suppressor --- capable models write somewhat longer responses,
partially masking the true capability signal in bivariate analysis. After residualization, the
capability signal dominates.

Figure~\ref{fig:bootstrap} (bootstrap distribution) confirms stability: CI $[0.976, 0.988]$
is narrow and far from zero.

\begin{table}[h]
\centering
\caption{H-E1 results summary}
\label{tab:he1}
\begin{tabular}{llll}
\toprule
Metric & Value & Threshold & Status \\
\midrule
$r_{\mathrm{partial}}$ & 0.9851 & $\geq 0.15$ & \textbf{EXCEEDED} \\
$p$-value & $1.69 \times 10^{-170}$ (one-tailed) & $< 0.05$ & \textbf{EXCEEDED} \\
Bootstrap CI lower & 0.976 & $> 0$ & \textbf{CONFIRMED} \\
VIF & 1.764 & $< 5.0$ & \textbf{MET} \\
\bottomrule
\end{tabular}
\end{table}

\subsection{RQ2: Capability Dominates Verbosity in OLS (H-M1)}

In standardized OLS:

\begin{quote}
$|\beta_{\mathrm{win}}| = 21.34$ vs $|\beta_{\mathrm{len}}| = 4.37$ --- dominance ratio
$4.88{:}1$ ($p_{\mathrm{win}} = 4.58 \times 10^{-145}$)
\end{quote}

Full model $R^2 = 0.963$ vs verbosity-only $R^2 = 0.256$ ($+70.7$pp from capability).

Figure~\ref{fig:coef} shows the standardized coefficient comparison.
$\beta_{\mathrm{len}} = -4.37$ (negative): verbosity is penalized once capability is
controlled --- confirming the LC correction is directionally correct. Mild heteroscedasticity
(BP $p = 0.011$) does not affect conclusion given $p_{\mathrm{win}} = 4.58 \times 10^{-145}$.

\begin{table}[h]
\centering
\caption{H-M1 standardized OLS results}
\label{tab:hm1}
\begin{tabular}{ll}
\toprule
Metric & Value \\
\midrule
$|\beta_{\mathrm{win\_std}}|$ & 21.34 \\
$|\beta_{\mathrm{len\_std}}|$ & $-4.37$ \\
Dominance ratio & $4.88\times$ \\
Full model $R^2$ & 0.963 \\
Verbosity-only $R^2$ & 0.256 \\
$R^2$ gain & $+70.7$pp \\
\bottomrule
\end{tabular}
\end{table}

\subsection{RQ3: Estimator Robustness via FWL-Inspired Residualization (H-M2)}

\begin{quote}
$\rho_{\mathrm{resid}} = 0.9739$, $p = 2.37 \times 10^{-144}$;
delta $= 0.0112$ ($< 0.02$ \checkmark)
\end{quote}

Two independent estimators converge within 1.1pp (Figure~\ref{fig:fwl}). While the FWL
theorem guarantees exact equality for Pearson OLS, our Spearman residualization provides an
independent approximation; the 1.1pp gap is well within sampling variation for $N=223$, ruling
out methodological artifact. Bootstrap CI on $\rho_{\mathrm{resid}}$: $[0.962, 0.981]$.

\begin{table}[h]
\centering
\caption{H-M2 FWL estimator comparison}
\label{tab:hm2}
\begin{tabular}{lll}
\toprule
Estimator & Value & $p$-value \\
\midrule
Pingouin partial\_corr & 0.9851 & $1.69 \times 10^{-170}$ (one-tailed) \\
FWL residual Spearman & 0.9739 & $2.37 \times 10^{-144}$ \\
Delta & 0.0112 & $< 0.02$ \checkmark \\
\bottomrule
\end{tabular}
\end{table}

\subsection{RQ4: LC Monotonicity Across Quartiles (H-C1)}

\begin{quote}
KW $H = 196.32$, $p = 2.63 \times 10^{-42}$, $\varepsilon^2 = 0.883$ (large);
Dunn Q1 vs Q4 $p = 1.04 \times 10^{-38}$

Monotonic: Q1 $= 7.14 <$ Q2 $= 14.69 <$ Q3 $= 26.41 <$ Q4 $= 51.62$
\end{quote}

Figure~\ref{fig:quartile_lc} shows the \texttt{LC\_winrate} boxplot by quartile.
$\varepsilon^2 = 0.883$ means ${\sim}88\%$ of \texttt{LC\_winrate} variance across models is
explained by capability quartile membership. Q4 models show $7.2\times$ higher median LC than
Q1. Bootstrap CI on Dunn $p$: $[1.43 \times 10^{-40}, 2.23 \times 10^{-35}]$.

\begin{table}[h]
\centering
\caption{H-C1 quartile results}
\label{tab:hc1}
\begin{tabular}{lll}
\toprule
Metric & Value & Status \\
\midrule
KW $p$ & $2.63 \times 10^{-42}$ & \textbf{PASS} \\
$\varepsilon^2$ & 0.883 & \textbf{LARGE} \\
Dunn Q1 vs Q4 $p$ & $1.04 \times 10^{-38}$ & \textbf{PASS} \\
Monotonic ordering & True & \textbf{CONFIRMED} \\
\bottomrule
\end{tabular}
\end{table}

\subsection{RQ5: Alignment Gap $\Delta$ and DV Choice (H-M3)}

Using $\Delta = \texttt{LC\_winrate} - \texttt{win\_rate}$ as DV:

\begin{quote}
KW $H = 22.19$, $p = 5.97 \times 10^{-5}$, $\varepsilon^2 = 0.088$ (medium);
Dunn Q1 vs Q4 $p = 1.0$ (n.s.)

Non-monotonic: Q1 $= 2.19$, Q2 $= 2.96$, Q3 $= 5.43$, Q4 $= 0.91$
\end{quote}

Figure~\ref{fig:delta} shows the $\Delta$ boxplot; Figure~\ref{fig:comparison} compares
$\varepsilon^2 = 0.088$ (H-M3) vs $\varepsilon^2 = 0.883$ (H-C1) --- a $10\times$ difference
from DV choice alone.

\begin{table}[h]
\centering
\caption{DV choice comparison: LC\_winrate vs $\Delta$}
\label{tab:dvchoice}
\begin{tabular}{lllll}
\toprule
DV & $\varepsilon^2$ & KW $p$ & Monotonic & Dunn Q1 vs Q4 \\
\midrule
\texttt{LC\_winrate} & 0.883 & $2.63 \times 10^{-42}$ & \checkmark & $p = 1.04 \times 10^{-38}$ \\
$\Delta$ & 0.088 & $5.97 \times 10^{-5}$ & $\times$ & $p = 1.0$ (n.s.) \\
\bottomrule
\end{tabular}
\end{table}

\subsection{Summary}

\begin{table}[h]
\centering
\caption{All hypothesis results}
\label{tab:summary}
\begin{tabular}{lllll}
\toprule
Hypothesis & Gate & Result & Key Metric \\
\midrule
H-E1 & MUST\_WORK & \textbf{PASS} & $r_{\mathrm{partial}} = 0.985 \gg 0.15$ \\
H-M1 & MUST\_WORK & \textbf{PASS} & $|\beta_{\mathrm{win}}| = 21.34 \gg |\beta_{\mathrm{len}}| = 4.37$ \\
H-M2 & SHOULD\_WORK & \textbf{PASS} & delta $= 0.011 < 0.02$ \\
H-M3 & SHOULD\_WORK & \textbf{PASS} & KW $p = 5.97 \times 10^{-5} < 0.05$ \\
H-C1 & SHOULD\_WORK & \textbf{PASS} & KW $p = 2.63 \times 10^{-42}$, Dunn $p = 1.04 \times 10^{-38}$ \\
\bottomrule
\end{tabular}
\end{table}
